% ============================================================
%   CHAPITRE : Géométrie dans l'espace
%   Niveau   : Terminale S / Terminale C
% ============================================================
\chapter{Géométrie dans l'espace}

\begin{citationchapitre}
	« La géométrie n'est pas vraie, elle est avantageuse. »\\[0.4em]
	— Henri Poincaré
\end{citationchapitre}

\begin{objectifs}
	\begin{itemize}[label=\textbullet]
		\item Repérer points et vecteurs de l'espace ; décomposer un vecteur.
		\item Calculer un \textbf{produit scalaire} et un \textbf{produit
		vectoriel} (formes analytique et géométrique) ; caractériser
		colinéarité, orthogonalité, alignement, coplanarité.
		\item Déterminer les équations \textbf{cartésiennes} et
		\textbf{paramétriques} d'une droite et d'un plan.
		\item Étudier les positions relatives de deux plans, d'un plan et d'une
		droite, de deux droites.
		\item Calculer des distances (point–plan, point–droite, entre deux
		plans) et construire un projeté orthogonal.
		\item Reconnaître l'équation d'une sphère et son intersection avec un plan.
	\end{itemize}
\end{objectifs}


% ============================================================
\section{Vecteurs et repérage dans l'espace}
% ============================================================

\begin{definition}
	L'espace est muni d'un \textbf{repère orthonormé direct}
	$(O\,;\,\vec\imath,\vec\jmath,\vec k)$.
	\begin{itemize}[label=\textbullet]
		\item Un point $M$ est repéré par ses coordonnées $M(x\,;\,y\,;\,z)$.
		\item Un vecteur $\vec u$ est repéré par ses coordonnées $\vec u(x\,;\,y\,;\,z)$,
		c'est-à-dire $\vec u = x\,\vec\imath + y\,\vec\jmath + z\,\vec k$.
	\end{itemize}
\end{definition}

\begin{propriete}{Calculs sur les coordonnées}{esp_coord}
	Soient $A(x_A;y_A;z_A)$, $B(x_B;y_B;z_B)$ et $\vec u(x;y;z)$, $\vec v(x';y';z')$.
	\begin{itemize}[label=\textbullet]
		\item $\vect{AB}\,(x_B - x_A\,;\,y_B - y_A\,;\,z_B - z_A)$.
		\item $\vec u + \vec v\,(x + x'\,;\,y + y'\,;\,z + z')$ \quad et \quad
		$\lambda\vec u\,(\lambda x\,;\,\lambda y\,;\,\lambda z)$.
		\item Norme : $\left\|\vec u\right\| = \sqrt{x^2 + y^2 + z^2}$ \quad ;
		\quad distance : $AB = \left\|\vect{AB}\right\|$.
		\item Milieu de $[AB]$ : $I\!\left(\dfrac{x_A+x_B}{2}\,;\,\dfrac{y_A+y_B}{2}\,;\,\dfrac{z_A+z_B}{2}\right)$.
	\end{itemize}
\end{propriete}

\begin{center}
\begin{tikzpicture}[>=Stealth, scale=1.0]
	% axes
	\draw[->] (0,0) -- (-1.5,-1.0) node[below left,font=\small] {$x$};
	\draw[->] (0,0) -- (3.4,0) node[right,font=\small] {$y$};
	\draw[->] (0,0) -- (0,3.0) node[above,font=\small] {$z$};
	\node[below right,font=\small] at (0,0) {$O$};
	% point M : projection sur (Oy) en (2.2,0), composante x = (-0.9,-0.6), z = (0,1.8)
	\coordinate (P) at (2.2,0);
	\coordinate (Q) at (1.3,-0.6);
	\coordinate (M) at (1.3,1.2);
	\draw[gray!55,dashed] (0,0) -- (P) -- (Q) -- (M);
	\draw[gray!55,dashed] (0,0) -- (Q);
	\draw[->,blue!70!black,thick] (0,0) -- (M);
	\filldraw[blue!70!black] (M) circle (1.6pt) node[above right,font=\small] {$M(x\,;y\,;z)$};
\end{tikzpicture}
\end{center}


% ============================================================
\section{Colinéarité, coplanarité, décomposition}
% ============================================================

\begin{propriete}{Colinéarité}{esp_colin}
	$\vec u$ et $\vec v$ ($\vec u \neq \vec 0$) sont \textbf{colinéaires} si et
	seulement s'il existe $k \in \R$ tel que $\vec v = k\,\vec u$, c'est-à-dire
	si leurs coordonnées sont \textbf{proportionnelles}.

	Trois points $A$, $B$, $C$ sont \textbf{alignés} $\iff$ $\vect{AB}$ et
	$\vect{AC}$ sont colinéaires.
\end{propriete}

\begin{propriete}{Coplanarité et décomposition}{esp_coplan}
	Soient $\vec u$ et $\vec v$ deux vecteurs \textbf{non colinéaires}.
	\begin{itemize}[label=\textbullet]
		\item $\vec w$ est \textbf{coplanaire} avec $\vec u$ et $\vec v$ $\iff$
		il existe $(a,b) \in \R^2$ tel que $\vec w = a\,\vec u + b\,\vec v$.
		\item Si $(\vec u,\vec v,\vec w)$ ne sont \textbf{pas} coplanaires, ils
		forment une \textbf{base} : tout vecteur de l'espace s'écrit de manière
		unique $a\,\vec u + b\,\vec v + c\,\vec w$.
	\end{itemize}
	Quatre points $A$, $B$, $C$, $D$ sont \textbf{coplanaires} $\iff$ il existe
	$(a,b) \in \R^2$ tel que $\vect{AD} = a\,\vect{AB} + b\,\vect{AC}$
	(avec $A$, $B$, $C$ non alignés).
\end{propriete}

\begin{astucebac}
	Le \textbf{choix d'une décomposition pertinente} est souvent la clé :
	pour un problème d'alignement ou de coplanarité, on exprime tous les
	vecteurs utiles dans \textbf{une même base} adaptée à la figure (par
	exemple $(\vect{AB}, \vect{AD}, \vect{AE})$ pour un cube $ABCDEFGH$).
\end{astucebac}


% ============================================================
\section{Produit scalaire}
% ============================================================

\begin{definition}[Deux expressions du produit scalaire]
	Pour $\vec u(x;y;z)$ et $\vec v(x';y';z')$ :
	$$
	\begin{aligned}
		\vec u \cdot \vec v &= x x' + y y' + z z'
		&&\text{(expression analytique)}\\
		\vec u \cdot \vec v &= \left\|\vec u\right\|\,\left\|\vec v\right\|\cos\theta
		&&\text{(expression géométrique, } \theta = (\vec u,\vec v) \text{)}
	\end{aligned}
	$$
\end{definition}

\begin{propriete}{Propriétés}{esp_ps}
	\begin{itemize}[label=\textbullet]
		\item $\vec u \cdot \vec u = \left\|\vec u\right\|^2$ ; symétrie
		$\vec u \cdot \vec v = \vec v \cdot \vec u$ ; bilinéarité.
		\item \textbf{Orthogonalité} : $\vec u \perp \vec v \iff \vec u \cdot \vec v = 0$.
		\item $\left\|\vec u + \vec v\right\|^2
		= \left\|\vec u\right\|^2 + 2\,\vec u\cdot\vec v + \left\|\vec v\right\|^2$.
		\item $\cos\theta = \dfrac{\vec u \cdot \vec v}{\left\|\vec u\right\|\,\left\|\vec v\right\|}$.
	\end{itemize}
\end{propriete}

\begin{exemple}
	$\vec u(1;-2;2)$ et $\vec v(3;0;-1)$ : $\vec u\cdot\vec v = 3 + 0 - 2 = 1$.
	$\left\|\vec u\right\| = 3$, $\left\|\vec v\right\| = \sqrt{10}$, donc
	$\cos\theta = \dfrac{1}{3\sqrt{10}}$.
\end{exemple}


% ============================================================
\section{Produit vectoriel}
% ============================================================

\begin{definition}[Produit vectoriel]
	Le \textbf{produit vectoriel} de $\vec u$ et $\vec v$, noté
	$\vec u \wedge \vec v$, est l'unique vecteur tel que :
	\begin{itemize}[label=\textbullet]
		\item $\vec u \wedge \vec v$ est \textbf{orthogonal} à $\vec u$ et à $\vec v$ ;
		\item $\left\|\vec u \wedge \vec v\right\| = \left\|\vec u\right\|\,\left\|\vec v\right\|\,\left|\sin\theta\right|$ ;
		\item $(\vec u,\ \vec v,\ \vec u \wedge \vec v)$ est une base \textbf{directe}.
	\end{itemize}

	\textbf{Expression analytique} : pour $\vec u(x;y;z)$, $\vec v(x';y';z')$,
	$$
	\vec u \wedge \vec v =
	\bigl(\,y z' - z y'\ ;\ z x' - x z'\ ;\ x y' - y x'\,\bigr).
	$$
\end{definition}

\begin{propriete}{Propriétés}{esp_pv}
	\begin{itemize}[label=\textbullet]
		\item $\vec v \wedge \vec u = -\,\vec u \wedge \vec v$ (antisymétrie) ;
		bilinéarité.
		\item $\vec u \wedge \vec v = \vec 0 \iff \vec u$ et $\vec v$ sont
		\textbf{colinéaires}.
		\item \textbf{Aire} du triangle $ABC$ :
		$\mathscr{A} = \dfrac{1}{2}\left\|\vect{AB} \wedge \vect{AC}\right\|$
		(l'aire du parallélogramme est $\left\|\vect{AB} \wedge \vect{AC}\right\|$).
		\item \textbf{Volume} du parallélépipède construit sur
		$\vec u,\vec v,\vec w$ : $V = \left|(\vec u \wedge \vec v)\cdot\vec w\right|$.
	\end{itemize}
\end{propriete}

\begin{exemple}
	$\vec u(1;2;3)$, $\vec v(0;1;-1)$ :
	$$
	\vec u \wedge \vec v =
	\bigl(2\times(-1) - 3\times1\ ;\ 3\times0 - 1\times(-1)\ ;\ 1\times1 - 2\times0\bigr)
	= (-5\ ;\ 1\ ;\ 1).
	$$
	On vérifie : $(-5)(1) + (1)(2) + (1)(3) = 0$ et $(-5)(0)+(1)(1)+(1)(-1) = 0$.
\end{exemple}


% ============================================================
\section{Le plan}
% ============================================================

\begin{propriete}{Plan défini par un point et deux vecteurs directeurs}{esp_plan_param}
	Le plan $(P)$ passant par $A$ et dirigé par $\vec u$, $\vec v$ (non
	colinéaires) admet la \textbf{représentation paramétrique}
	$$
	M \in (P) \iff
	\begin{cases}
		x = x_A + s\,x_{\vec u} + t\,x_{\vec v} \\
		y = y_A + s\,y_{\vec u} + t\,y_{\vec v} \\
		z = z_A + s\,z_{\vec u} + t\,z_{\vec v}
	\end{cases}
	\quad (s,t) \in \R^2.
	$$
\end{propriete}

\begin{propriete}{Plan défini par un point et un vecteur normal}{esp_plan_cart}
	Un vecteur $\vec n \neq \vec 0$ est \textbf{normal} à $(P)$ s'il est
	orthogonal à tout vecteur de $(P)$.
	$$
	M \in (P) \iff \vec n \cdot \vect{AM} = 0.
	$$
	Il en résulte une \textbf{équation cartésienne} :
	$$
	(P) : ax + by + cz + d = 0,
	\qquad \text{avec } \vec n\,(a\,;\,b\,;\,c) \text{ normal à } (P).
	$$
	Réciproquement, tout plan admet une équation de cette forme, et $\vec n(a;b;c)$
	en est un vecteur normal.
\end{propriete}

\begin{methode}
	\textbf{Équation cartésienne du plan $(ABC)$.}
	\begin{enumerate}
		\item Calculer $\vec n = \vect{AB} \wedge \vect{AC}$ (vecteur normal).
		\item L'équation est $a x + b y + c z + d = 0$ avec $\vec n(a;b;c)$.
		\item Déterminer $d$ en écrivant que $A \in (P)$.
	\end{enumerate}
\end{methode}

\begin{center}
\begin{tikzpicture}[>=Stealth, scale=1.0]
	% plan : parallélogramme
	\draw[fill=gray!12,draw=black!55] (-1.8,-0.9) -- (2.2,-0.5) -- (2.6,0.9) -- (-1.4,0.5) -- cycle;
	\node[font=\small] at (2.2,-1.0) {$(P)$};
	\coordinate (A) at (0.2,-0.1);
	\filldraw (A) circle (1.6pt) node[below,font=\small] {$A$};
	% vecteur normal
	\draw[->,red,thick] (A) -- ++(0.35,1.7) node[above,font=\small,red] {$\vec n$};
	% deux directeurs
	\draw[->,blue!70!black] (A) -- ++(1.4,0.15) node[right,font=\scriptsize,blue!70!black] {$\vec u$};
	\draw[->,blue!70!black] (A) -- ++(-0.5,0.55) node[left,font=\scriptsize,blue!70!black] {$\vec v$};
\end{tikzpicture}
\end{center}


% ============================================================
\section{La droite}
% ============================================================

\begin{propriete}{Représentation paramétrique d'une droite}{esp_droite_param}
	La droite $(D)$ passant par $A$ et de \textbf{vecteur directeur}
	$\vec u(\alpha;\beta;\gamma)$ est
	$$
	M \in (D) \iff
	\begin{cases}
		x = x_A + t\,\alpha \\
		y = y_A + t\,\beta \\
		z = z_A + t\,\gamma
	\end{cases}
	\quad t \in \R.
	$$
\end{propriete}

\remarque
Une droite de l'espace est l'\textbf{intersection de deux plans} ; elle peut
donc aussi être décrite par un \textbf{système de deux équations
cartésiennes}. Un vecteur directeur est alors $\vec n_1 \wedge \vec n_2$.


% ============================================================
\section{Positions relatives}
% ============================================================

\begin{propriete}{Deux plans}{esp_pos_pp}
	Soient $(P) : \vec n$ et $(P') : \vec n'$.
	\begin{itemize}[label=\textbullet]
		\item $\vec n$ et $\vec n'$ \textbf{colinéaires} : $(P)$ et $(P')$ sont
		\textbf{parallèles} (confondus ou strictement parallèles).
		\item sinon : $(P)$ et $(P')$ sont \textbf{sécants} suivant une droite de
		vecteur directeur $\vec n \wedge \vec n'$.
	\end{itemize}
\end{propriete}

\begin{propriete}{Un plan et une droite}{esp_pos_pd}
	Soient $(D) : \vec u$ et $(P) : \vec n$.
	\begin{itemize}[label=\textbullet]
		\item $\vec u \cdot \vec n = 0$ : $(D)$ est \textbf{parallèle} à $(P)$
		($(D) \subset (P)$ ou strictement parallèle — on teste avec un point).
		\item $\vec u \cdot \vec n \neq 0$ : $(D)$ et $(P)$ sont \textbf{sécants}
		en un point (on remplace la paramétrisation de $(D)$ dans l'équation de $(P)$).
		\item $\vec u$ et $\vec n$ \textbf{colinéaires} : $(D)$ est
		\textbf{perpendiculaire} à $(P)$.
	\end{itemize}
\end{propriete}

\begin{propriete}{Deux droites}{esp_pos_dd}
	Soient $(D) : A, \vec u$ et $(D') : A', \vec u'$.
	\begin{itemize}[label=\textbullet]
		\item Si $\bigl(\vec u,\ \vec u',\ \vect{AA'}\bigr)$ sont
		\textbf{coplanaires} (produit mixte nul), $(D)$ et $(D')$ sont
		\textbf{coplanaires} : parallèles si $\vec u,\vec u'$ colinéaires,
		sécantes sinon.
		\item Sinon, $(D)$ et $(D')$ sont \textbf{non coplanaires}.
		\item $\vec u \cdot \vec u' = 0$ : $(D)$ et $(D')$ sont
		\textbf{orthogonales} (sans être nécessairement sécantes).
	\end{itemize}
\end{propriete}

\begin{center}
\renewcommand{\arraystretch}{1.3}
\begin{tabular}{>{\raggedright\arraybackslash}p{3.0cm}
                >{\raggedright\arraybackslash}p{2.3cm}
                >{\raggedright\arraybackslash}p{2.6cm}}
\hline
\textbf{Objets} & \textbf{Test} & \textbf{Conclusion} \\
\hline
Plan / plan & $\vec n \parallel \vec n'$ & parallèles \\
 & sinon & sécants (droite) \\
Droite / plan & $\vec u \cdot \vec n = 0$ & parallèles \\
 & sinon & sécants (point) \\
Droite / droite & mixte $= 0$ & coplanaires \\
 & sinon & non coplanaires \\
\hline
\end{tabular}
\end{center}


% ============================================================
\section{Distances et projeté orthogonal}
% ============================================================

\begin{propriete}{Distance d'un point à un plan}{esp_dist_pp}
	Pour $(P) : ax + by + cz + d = 0$ et $M(x_M;y_M;z_M)$ :
	$$
	d\bigl(M,(P)\bigr) = \frac{\left|a x_M + b y_M + c z_M + d\right|}{\sqrt{a^2 + b^2 + c^2}}.
	$$
	La \textbf{distance entre deux plans parallèles} est la distance d'un point
	de l'un à l'autre.
\end{propriete}

\begin{propriete}{Distance d'un point à une droite}{esp_dist_pd}
	Pour $(D)$ passant par $B$ de vecteur directeur $\vec u$, et un point $A$ :
	$$
	d\bigl(A,(D)\bigr) = \frac{\left\|\vect{BA} \wedge \vec u\right\|}{\left\|\vec u\right\|}.
	$$
\end{propriete}

\begin{methode}
	\textbf{Projeté orthogonal.}
	\begin{itemize}[label=\textbullet]
		\item de $M$ sur $(P)$ : point d'intersection de $(P)$ avec la droite
		passant par $M$ de vecteur directeur $\vec n$ (normal à $(P)$) ;
		\item de $A$ sur $(D)$ : point $H \in (D)$ tel que $\vect{AH} \cdot \vec u = 0$.
	\end{itemize}
	La distance cherchée est alors $MH$ (resp. $AH$).
\end{methode}


% ============================================================
\section{La sphère}
% ============================================================

\begin{propriete}{Équation d'une sphère}{esp_sphere}
	La sphère de centre $\Omega(a;b;c)$ et de rayon $R$ a pour équation
	$$
	(x - a)^2 + (y - b)^2 + (z - c)^2 = R^2.
	$$
	Une équation de la forme
	$x^2 + y^2 + z^2 - 2a x - 2b y - 2c z + \delta = 0$ est celle d'une sphère
	de centre $(a;b;c)$ et de rayon $R = \sqrt{a^2 + b^2 + c^2 - \delta}$
	(si ce nombre est strictement positif).
\end{propriete}

\begin{propriete}{Intersection sphère–plan}{esp_sphere_plan}
	Soit $\Omega$ le centre, $R$ le rayon, et $h = d\bigl(\Omega,(P)\bigr)$.
	\begin{itemize}[label=\textbullet]
		\item $h < R$ : l'intersection est un \textbf{cercle} de rayon
		$\sqrt{R^2 - h^2}$, de centre le projeté orthogonal de $\Omega$ sur $(P)$ ;
		\item $h = R$ : $(P)$ est \textbf{tangent} (un seul point) ;
		\item $h > R$ : l'intersection est \textbf{vide}.
	\end{itemize}
\end{propriete}


% ============================================================
\section{Méthodes et exercices résolus}
% ============================================================

\exerciceresolu
\textbf{Plan $(ABC)$ et position d'une droite}

On donne $A(1;0;2)$, $B(2;-1;0)$, $C(0;1;1)$ et la droite $(\Delta)$ de
représentation paramétrique
$\begin{cases} x = 1 + 2t \\ y = t \\ z = 3 - t \end{cases}$, $t \in \R$.
\begin{enumerate}
	\item Déterminer une équation cartésienne du plan $(ABC)$.
	\item Étudier la position relative de $(\Delta)$ et de $(ABC)$.
\end{enumerate}

\solution

\textbf{1.} $\vect{AB}(1;-1;-2)$, $\vect{AC}(-1;1;-1)$.
$$
\vec n = \vect{AB} \wedge \vect{AC} =
\begin{pmatrix}
	(-1)(-1) - (-2)(1) \\
	(-2)(-1) - (1)(-1) \\
	(1)(1) - (-1)(-1)
\end{pmatrix}
= \begin{pmatrix} 3 \\ 3 \\ 0 \end{pmatrix}.
$$
On prend $\vec n(1;1;0)$ : $(ABC) : x + y + d = 0$. Avec $A$ : $1 + 0 + d = 0$,
donc $d = -1$.
$$
\boxed{(ABC) : x + y - 1 = 0.}
$$

\textbf{2.} Vecteur directeur de $(\Delta)$ : $\vec u(2;1;-1)$.
$\vec u \cdot \vec n = 2 + 1 + 0 = 3 \neq 0$ : $(\Delta)$ et $(ABC)$ sont
\textbf{sécants} en un point. En reportant : $(1 + 2t) + t - 1 = 0$, soit
$3t = 0$, $t = 0$. Le point d'intersection est $I(1;0;3)$.

\exerciceresolu
\textbf{Distance et projeté orthogonal}

On donne le point $A(3;1;-1)$ et le plan $(P) : 2x - y + 2z - 3 = 0$.
Déterminer la distance de $A$ à $(P)$ et les coordonnées du projeté orthogonal
$H$ de $A$ sur $(P)$.

\solution

\textbf{Distance.}
$$
d\bigl(A,(P)\bigr) = \frac{\left|2(3) - 1 + 2(-1) - 3\right|}{\sqrt{4 + 1 + 4}}
= \frac{\left|6 - 1 - 2 - 3\right|}{3} = \frac{0}{3} = 0.
$$
Donc $A \in (P)$ et $H = A$.

\medskip
\emph{Variante : si on remplace $A$ par $A'(3;1;1)$ :}
$d\bigl(A',(P)\bigr) = \dfrac{|6 - 1 + 2 - 3|}{3} = \dfrac{4}{3}$.
La droite par $A'$ dirigée par $\vec n(2;-1;2)$ est
$(x;y;z) = (3 + 2t\,;\,1 - t\,;\,1 + 2t)$. En reportant dans $(P)$ :
$2(3+2t) - (1 - t) + 2(1+2t) - 3 = 0 \Rightarrow 4 + 9t = 0 \Rightarrow t = -\dfrac{4}{9}$.
D'où $H\!\left(\dfrac{19}{9}\,;\,\dfrac{13}{9}\,;\,\dfrac{1}{9}\right)$
et $A'H = \dfrac{4}{3}$.

\exerciceresolu
\textbf{Sphère et plan}

On considère la sphère $(S) : x^2 + y^2 + z^2 - 2x + 4y - 6z + 5 = 0$ et le
plan $(P) : x - 2y + 2z = 1$.
\begin{enumerate}
	\item Déterminer le centre $\Omega$ et le rayon $R$ de $(S)$.
	\item Étudier l'intersection de $(S)$ et de $(P)$.
\end{enumerate}

\solution

\textbf{1.} $x^2 - 2x = (x-1)^2 - 1$, $y^2 + 4y = (y+2)^2 - 4$,
$z^2 - 6z = (z-3)^2 - 9$. L'équation devient
$(x-1)^2 + (y+2)^2 + (z-3)^2 = 9$.
Donc $\Omega(1;-2;3)$ et $R = 3$.

\textbf{2.} $d\bigl(\Omega,(P)\bigr)
= \dfrac{\left|1 - 2(-2) + 2(3) - 1\right|}{\sqrt{1 + 4 + 4}}
= \dfrac{\left|1 + 4 + 6 - 1\right|}{3} = \dfrac{10}{3} > 3 = R$.
L'intersection est \textbf{vide}.

\exerciceresolu
\textbf{Deux droites : coplanaires ou non ?}

$(D) : (x;y;z) = (1 + t\,;\,2t\,;\,-t)$ et
$(D') : (x;y;z) = (2 + s\,;\,1 - s\,;\,s)$.
\begin{enumerate}
	\item Ces droites sont-elles parallèles ?
	\item Sont-elles coplanaires ?
\end{enumerate}

\solution

\textbf{1.} $\vec u(1;2;-1)$, $\vec u'(1;-1;1)$ : non proportionnels, donc
non parallèles.

\textbf{2.} $A(1;0;0) \in (D)$, $A'(2;1;0) \in (D')$, $\vect{AA'}(1;1;0)$.
Produit mixte :
$$
(\vec u \wedge \vec u') \cdot \vect{AA'}.
$$
$\vec u \wedge \vec u' = (2\cdot1 - (-1)(-1)\ ;\ (-1)(1) - 1\cdot1\ ;\ 1\cdot(-1) - 2\cdot1)
= (1\ ;\ -2\ ;\ -3)$.
$(\vec u \wedge \vec u') \cdot \vect{AA'} = 1 - 2 + 0 = -1 \neq 0$.
Les droites sont \textbf{non coplanaires}.


% ============================================================
\section{Exercices d'entraînement}
% ============================================================

\subsection{Échauffement : lire une figure}

\exercice{Le cube \hfill \facile}
$ABCDEFGH$ est un cube d'arête $1$. On munit l'espace du repère orthonormé
$\left(A\,;\,\vect{AB},\vect{AD},\vect{AE}\right)$.

\begin{center}
\begin{tikzpicture}[>=Stealth, scale=1.7]
	\coordinate (A) at (0,0);
	\coordinate (B) at (1,0);
	\coordinate (C) at (1.5,0.5);
	\coordinate (D) at (0.5,0.5);
	\coordinate (E) at (0,1);
	\coordinate (F) at (1,1);
	\coordinate (G) at (1.5,1.5);
	\coordinate (H) at (0.5,1.5);
	% arêtes visibles
	\draw[thick] (A)--(B)--(C)--(G)--(H)--(E)--(A);
	\draw[thick] (B)--(F)--(G) (E)--(F);
	% arêtes cachées (issues du sommet D)
	\draw[thick,dashed] (D)--(A) (D)--(C) (D)--(H);
	\foreach \p/\pos in {A/below left,B/below right,C/right,D/left,E/left,F/above right,G/right,H/above}
	  \filldraw (\p) circle (0.6pt) node[\pos,font=\scriptsize] {$\p$};
\end{tikzpicture}
\end{center}

\begin{enumerate}
	\item Donner les coordonnées des huit sommets.
	\item Calculer les coordonnées du centre $I$ du cube et du milieu $J$ de $[FG]$.
	\item Calculer $\vect{AG} \cdot \vect{BH}$. Les droites $(AG)$ et $(BH)$
	sont-elles orthogonales ?
	\item Les points $A$, $C$, $F$ sont-ils alignés ? Le point $H$ est-il dans
	le plan $(ACF)$ ?
\end{enumerate}

\exercice{Le tétraèdre \hfill \moyen}
$ABCD$ est un tétraèdre. On donne $A(0;0;0)$, $B(2;0;0)$, $C(0;2;0)$,
$D(0;0;2)$.
\begin{enumerate}
	\item Calculer $\vect{AB} \wedge \vect{AC}$, puis l'aire du triangle $ABC$.
	\item En déduire une équation cartésienne du plan $(ABC)$, puis du plan $(BCD)$.
	\item Calculer la distance de $A$ au plan $(BCD)$, puis le volume du
	tétraèdre.
\end{enumerate}

\exercice{Vrai ou faux \hfill \facile}
Répondre par \emph{vrai} ou \emph{faux} en justifiant brièvement.
\begin{enumerate}
	\item Si $\vec u \cdot \vec v = 0$ alors $\vec u = \vec 0$ ou $\vec v = \vec 0$.
	\item $\vec u \wedge \vec v$ est orthogonal à la fois à $\vec u$ et à $\vec v$.
	\item Deux droites orthogonales de l'espace sont sécantes.
	\item Le vecteur $\vec n(a;b;c)$ est normal au plan d'équation $ax+by+cz+d=0$.
	\item Si deux plans ont des vecteurs normaux colinéaires, ils sont confondus.
	\item Trois points de l'espace sont toujours coplanaires.
	\item $\left\|\vec u \wedge \vec v\right\|$ est l'aire du parallélogramme
	construit sur $\vec u$ et $\vec v$.
\end{enumerate}

\subsection{Vecteurs, produit scalaire et produit vectoriel}

\exercice{Décomposition et colinéarité \hfill \moyen}
On donne $\vec u(1;2;-1)$, $\vec v(2;-1;3)$, $\vec w(4;3;1)$.
\begin{enumerate}
	\item Montrer que $\vec u$ et $\vec v$ ne sont pas colinéaires.
	\item Le vecteur $\vec w$ est-il coplanaire avec $\vec u$ et $\vec v$ ?
	Si oui, écrire $\vec w = a\,\vec u + b\,\vec v$.
	\item Déterminer un vecteur orthogonal à la fois à $\vec u$ et à $\vec v$.
\end{enumerate}

\exercice{Produit scalaire — les deux expressions \hfill \moyen}
$A(1;1;0)$, $B(3;0;2)$, $C(2;3;1)$.
\begin{enumerate}
	\item Calculer $\vect{AB} \cdot \vect{AC}$ de deux manières
	(analytique, puis via les longueurs et l'angle).
	\item En déduire une mesure de l'angle $\widehat{BAC}$ (au degré près).
	\item Le triangle $ABC$ est-il rectangle ?
\end{enumerate}

\exercice{Produit vectoriel et aire \hfill \moyen}
$A(1;0;1)$, $B(0;2;3)$, $C(2;1;0)$.
\begin{enumerate}
	\item Calculer $\vect{AB} \wedge \vect{AC}$.
	\item En déduire l'aire du triangle $ABC$.
	\item Déterminer une équation cartésienne du plan $(ABC)$.
\end{enumerate}

\subsection{Droites et plans}

\exercice{Équations d'un plan \hfill \moyen}
\begin{enumerate}
	\item Écrire une équation cartésienne du plan passant par $A(2;-1;3)$ de
	vecteur normal $\vec n(1;2;-2)$.
	\item Écrire une représentation paramétrique du plan passant par
	$B(1;0;0)$ dirigé par $\vec u(1;1;0)$ et $\vec v(0;1;1)$.
	\item Ces deux plans sont-ils parallèles ? sécants ?
\end{enumerate}

\exercice{Intersection droite–plan \hfill \moyen}
Soit $(P) : x + 2y - z + 1 = 0$ et $(D)$ passant par $A(1;1;2)$ de vecteur
directeur $\vec u(2;-1;1)$.
\begin{enumerate}
	\item Montrer que $(D)$ et $(P)$ sont sécants.
	\item Déterminer les coordonnées du point d'intersection.
	\item Déterminer l'équation du plan contenant $A$ et perpendiculaire à $(D)$.
\end{enumerate}

\exercice{Positions relatives de deux droites \hfill \difficile}
$(D) : (x;y;z) = (1 + 2t\,;\,-t\,;\,3 + t)$ et
$(D') : (x;y;z) = (s\,;\,1 + s\,;\,2 - s)$.
\begin{enumerate}
	\item $(D)$ et $(D')$ sont-elles parallèles ?
	\item Sont-elles coplanaires ? Si oui, sont-elles sécantes ? Préciser le
	point ou le plan qui les contient.
	\item Sont-elles orthogonales ?
\end{enumerate}

\subsection{Sujets types Bac}

\exercice{Sujet type Bac — plan, droite et distance \hfill \difficile}
L'espace est muni d'un repère orthonormé direct
$(O\,;\,\vec\imath,\vec\jmath,\vec k)$. On donne
$A(1;-1;1)$, $B(3;2;1)$ et $C(2;-1;0)$.
\begin{enumerate}
	\item Vérifier que le plan $(ABC)$ a pour équation $3x - 2y + 3z - 8 = 0$. \hfill (0,25 pt)
	\item Écrire une équation cartésienne du plan $(Q)$ contenant $B$, de
	vecteur normal $\vect{AC}$. \hfill (0,5 pt)
	\item
	\begin{enumerate}
		\item Montrer que $(ABC) \cap (Q)$ est une droite $(D)$ dont on donnera
		une représentation paramétrique. \hfill (0,75 pt)
		\item Montrer que $C$ est le projeté orthogonal de $A$ sur $(D)$ et en
		déduire la distance de $A$ à $(D)$. \hfill (0,75 pt)
	\end{enumerate}
	\item Soit $(R) : x - y + z = 1$. Déterminer les coordonnées du point $I$,
	intersection des trois plans $(ABC)$, $(Q)$ et $(R)$. \hfill (0,75 pt)
\end{enumerate}

\exercice{Sujet type Bac — produit vectoriel et parallélisme \hfill \difficile}
On donne $A(2;3;1)$, $B(-1;1;1)$, $C(-3;1;0)$, $D(4;-1;3)$.
\begin{enumerate}
	\item Calculer $\vec n = \vect{AB} \wedge \vect{AC}$. \hfill (0,5 pt)
	\item En déduire une équation cartésienne du plan $(ABC)$. \hfill (0,5 pt)
	\item Donner une représentation paramétrique de la droite $(\Delta)$ de
	vecteur directeur $\vec u(1;2;-1)$ passant par $D$. \hfill (0,5 pt)
	\item Montrer que $(\Delta)$ est strictement parallèle au plan $(ABC)$. \hfill (0,5 pt)
\end{enumerate}

\exercice{Sujet type Bac — deux plans parallèles et orthogonalité \hfill \difficile}
On donne le point $A(2;2;1)$ et le plan $(P) : 3x + 2y + 6z + 33 = 0$.
\begin{enumerate}
	\item Vérifier que $(P)$ est strictement parallèle au plan $(Q)$ contenant
	$A$, de vecteur normal $\vec n(-3;-2;-6)$. \hfill (0,25 pt)
	\item Déterminer la distance qui sépare $(P)$ et $(Q)$. \hfill (0,75 pt)
	\item Soit $H(-1;0;-5)$. Justifier que $H \in (P)$ et que $(AH)$ et $(P)$
	sont orthogonaux. \hfill (0,75 pt)
	\item Calculer $\left\|\vect{AH}\right\|$ et retrouver la distance de $A$ à
	$(P)$. \hfill (0,5 pt)
\end{enumerate}

\exercice{Sujet type Bac — sphère et droite perpendiculaire à un plan \hfill \difficile}
On donne $A(1;-1;1)$, $B(-2;1;-1)$, $C(3;2;2)$.
\begin{enumerate}
	\item
	\begin{enumerate}
		\item Calculer $\vect{AB} \wedge \vect{AC}$. \hfill (0,25 pt)
		\item En déduire une équation cartésienne du plan $(ABC)$. \hfill (0,5 pt)
	\end{enumerate}
	\item Soit $(S) : x^2 + y^2 + z^2 - 4x + 4y - 2z + 8 = 0$. Déterminer le
	centre $\Omega$ et le rayon $R$ de $(S)$. \hfill (0,5 pt)
	\item Soit $(\Delta)$ la droite passant par $\Omega$ et perpendiculaire à
	$(ABC)$.
	\begin{enumerate}
		\item Donner une représentation paramétrique de $(\Delta)$. \hfill (0,25 pt)
		\item Déterminer les coordonnées du point $H$, intersection de
		$(\Delta)$ et de $(ABC)$. \hfill (0,5 pt)
		\item $(ABC)$ coupe-t-il la sphère $(S)$ ? Si oui, préciser le rayon du
		cercle d'intersection. \hfill (0,5 pt)
	\end{enumerate}
\end{enumerate}

\exercice{Sujet type Bac — deux droites orthogonales et plan commun \hfill \difficile}
On donne le point $A(2;-1;3)$ et le vecteur $\vec u(3;-2;1)$.
\begin{enumerate}
	\item Écrire une représentation paramétrique de la droite $(D)$ passant par
	$A$ de vecteur directeur $\vec u$. \hfill (0,25 pt)
	\item Soit $(D') : (x;y;z) = (4t + 1\,;\,5t - 8\,;\,-2t + 6)$, $t \in \R$.
	\begin{enumerate}
		\item Montrer que $(D)$ et $(D')$ sont orthogonales. \hfill (0,75 pt)
		\item Déterminer les coordonnées du point $I$, intersection de $(D)$ et
		$(D')$. \hfill (0,5 pt)
	\end{enumerate}
	\item Écrire une équation cartésienne du plan $(P)$ contenant $(D)$ et
	$(D')$. \hfill (0,75 pt)
\end{enumerate}
